\documentclass[10pt,a4paper]{amsart}
\usepackage[margin=19mm]{geometry}
\usepackage{amsmath,amssymb,mathtools}
\usepackage[scr=rsfs,cal=euler]{mathalfa}
\usepackage{xfrac}
\usepackage[unicode,hidelinks,bookmarks=false]{hyperref}
\newcommand{\ie}{i.e.\ }
\newcommand{\cf}{cf.\ }
\newcommand{\resp}{respectively\ }
\newcommand{\Subsection}{\subsection}
\providecommand{\nobreakdash}{\nobreak\hbox{-}}
\setlength{\emergencystretch}{2em}
\multlinegap0pt
\usepackage{amssymb}
\newtheorem{thm}{Theorem}
\newtheorem{lem}{Lemma}
\newtheorem{prp}{Proposition}
\newcommand{\C}{\mathbb{C}}
\newcommand{\N}{\mathbb{N}}
\newcommand{\R}{\mathbb{R}}
\newcommand{\dist}{\operatorname{dist}}
\newcommand{\inrad}{\operatorname{inrad}}
\newcommand{\vol}{\operatorname{vol}}
\title{On the inner radius of the nonvanishing set for eigenfunctions of complex elliptic operators}
\author{Omer Friedland, Henrik Uebersch\"ar}
\date{Source: arXiv:2603.10602v1 (CC BY 4.0)}
\begin{document}
\maketitle

\noindent\emph{Source and licence.} On the inner radius of the nonvanishing set for eigenfunctions of complex elliptic operators. Version arXiv:2603.10602v1 (CC BY 4.0), distributed by arXiv under the Creative Commons Attribution 4.0 International licence, http://creativecommons.org/licenses/by/4.0/, as recorded in the arXiv licence metadata for this exact version and shown on the abstract page https://arxiv.org/abs/2603.10602v1. Full licence notice: \texttt{licenses/arxiv-2603.10602-CC-BY-4.0.txt}.\par\smallskip

\noindent\textit{Editorial scope.} Counted unit: the article's thm thm:main (source0.tex lines 123--130). The article's complete original proof of it is delivered with it (source0.tex lines 354--444, one \texttt{proof} environment of 91 lines, which the article places away from the statement). No mathematics is written, repaired or re-derived: the statement and the proof are the article's own lines, in the article's own notation, and no retained line is substituted or re-flowed.  Retained uncounted as the source-local context the proof needs: lines 85--87; lines 90--92; lines 166--178; lines 193--198; lines 264--277; lines 293--300; lines 316--327.  Nothing was omitted from any retained span, so the package carries no omission record.  No retained line contains a citation, so the wrapper carries no bibliography and no bibliography entry.  No retained line contains a figure environment, an image-inclusion command or drawing code, so no image is delivered and the package declares no asset dependency.  Editorial formatting only, in addition: an AMS \texttt{amsart} preamble that restates the article's own declarations and macros used by the retained lines, namely the environments \texttt{thm}, \texttt{lem}, \texttt{prp} on separate counters, together with the standard packages those lines need.  The retained environments print in this document as Theorem 1, Lemma 1, Prp 1 in order of appearance, whereas the article prints the same items with the numbering of its own full document.  The complete native source of the article as distributed in the arXiv e-print for this exact version is bundled unchanged as \texttt{sources/196149/source0.tex}.
\begin{equation}\label{eq:ellipticity}
|P(\xi)|\ \ge\ c_{\rm ell} |\xi|^m,\quad \xi\in\R^d.
\end{equation}

\begin{equation}\label{eq:eigen}
H\psi_\lambda = \lambda\psi_\lambda\quad\text{in }\Omega.
\end{equation}

\begin{lem}[Lipschitz nonvanishing ball]\label{lem:lipschitz-ball}
Let $\Omega\subset\R^d$ be open and let $G:\Omega\to\C$ be Lipschitz on $\Omega$ with Lipschitz constant $L$, i.e.
$|G(x)-G(y)|\le L|x-y|$ for all $x,y\in\Omega$.
Fix $x_0\in\Omega$ and assume $|G(x_0)|\ge \eta>0$.
Set
$$
\rho := \min\Bigl\{\frac{\eta}{2L},\ \dist(x_0,\Omega^c)\Bigr\}.
$$
Then
$$
B(x_0,\rho)\ \subset\ \{x\in\Omega:\ G(x)\neq 0\}.
$$
\end{lem}

\begin{lem}[Pointwise lower bound from $L^2$ mass]\label{lem:L2-to-Linfty}
Let $u\in L^2(B(0,R))$. Then
$$
\sup_{B(0,R)} |u|\ \ge\ \frac{\|u\|_{L^2(B(0,R))}}{\vol(B(0,R))^{1/2}}.
$$
\end{lem}

\begin{prp}[Uniform local Lipschitz bound]\label{prp:uniform-lip}
Let $H$ be as in Theorem~\ref{thm:main}.
Then there exists a constant $L_{d,H}>0$ such that for every $\mu\in\C$ with $1\le|\mu|\le 2$
and every solution $u\in L^2(B(0,1))$ of
$$
Hu = \mu u\quad\text{in }B(0,1),
$$
we have
$$
\sup_{x\in B(0,3/4)}|\nabla u(x)|\ \le\ L_{d,H} \|u\|_{L^2(B(0,1))}.
$$
In particular, $u$ is Lipschitz on $B(0,3/4)$ with Lipschitz constant at most
$L_{d,H}\|u\|_{L^2(B(0,1))}$.
\end{prp}

\begin{lem}[Bounded overlap cover]\label{lem:cover}
Let $E\subset\R^d$ and let $r>0$.
There exists a (finite or countable) set $\{x_j\}_{j\in J}\subset E$ with the following properties:
\begin{enumerate}
\item $E\subset \bigcup_{j\in J} B(x_j,r/2)$.
\item The family of dilated balls $\{B(x_j,r)\}_{j\in J}$ has bounded overlap: there exists $N_d\in\N$ depending only on $d$ such that each point of $\R^d$ belongs to at most $N_d$ balls $B(x_j,r)$.
\end{enumerate}
\end{lem}

\begin{lem}[Existence of a good ball]\label{lem:good-ball}
Let $r>0$ and set $E := \Omega_{-r}$.
Let $f\in L^1(\Omega)$ satisfy $f\ge 0$ and $\int_\Omega f>0$.
Define $M := \int_E f$.
Let $\{x_j\}_{j\in J}\subset E$ be as in Lemma~\ref{lem:cover}, so that $E\subset \bigcup_{j} B(x_j,r/2)$ and the balls $B(x_j,r)$ have overlap bounded by $N_d$.
Then there exists $j_0\in J$ such that
\begin{equation}\label{eq:ratio}
\frac{\int_{B(x_{j_0},r/2)} f}{\int_{B(x_{j_0},r)} f}
\ \ge\
\frac{M}{2N_d \int_\Omega f}.
\end{equation}
\end{lem}

\begin{thm}[Quantitative inradius bound]\label{thm:main}
Let $\Omega\subset\R^d$ be open and let $H$ be a homogeneous constant--coefficient elliptic operator of order $m$ satisfying \eqref{eq:ellipticity}.
Then there exists a constant $c_{d,H}>0$ (depending only on $d$ and $H$) such that for every $|\lambda|\ge 1$ and every nonzero solution $\psi_\lambda\in L^2(\Omega)$ of \eqref{eq:eigen} we have
\begin{equation}\label{eq:main}
\inrad(\Sigma_\lambda)\ \ge\ c_{d,H} r_\lambda 
\frac{\|\psi_\lambda\|_{L^2(\Omega_{-r_\lambda})}}{\|\psi_\lambda\|_{L^2(\Omega)}}.
\end{equation}
\end{thm}

\begin{proof}[Proof of Theorem~\ref{thm:main}]
Let $\psi_\lambda\in L^2(\Omega)$ be a nonzero solution of \eqref{eq:eigen}.
Set $r := r_\lambda = |\lambda|^{-1/m}$ and define
$$
M := \|\psi_\lambda\|_{L^2(\Omega_{-r})}^2,\quad
N := \|\psi_\lambda\|_{L^2(\Omega)}^2>0.
$$
If $M = 0$ then \eqref{eq:main} is trivial. Assume $M>0$.

\textbf{Step 1: choose a good ball at scale $r$.}
Apply Lemma~\ref{lem:good-ball} with $f = |\psi_\lambda|^2$, $E = \Omega_{-r}$.
Since $\int_\Omega f = N$, there exists $x_0\in\Omega_{-r}$ such that
\begin{equation}\label{eq:good-ratio}
\frac{\|\psi_\lambda\|_{L^2(B(x_0,r/2))}^2}{\|\psi_\lambda\|_{L^2(B(x_0,r))}^2}
\ \ge\ \frac{M}{2N_d N}.
\end{equation}
Because $x_0\in\Omega_{-r}$, we have $B(x_0,r)\subset\Omega$.

\textbf{Step 2: rescale to unit scale.}
Define the rescaled and normalized function $u:B(0,1)\to\C$ by
\begin{equation}\label{eq:rescale}
u(y) := \frac{r^{d/2} \psi_\lambda(x_0+r y)}{\|\psi_\lambda\|_{L^2(B(x_0,r))}}.
\end{equation}
Then $\|u\|_{L^2(B(0,1))} = 1$, and similarly
\begin{equation}\label{eq:u-half}
\|u\|_{L^2(B(0,1/2))}^2
 = \frac{\|\psi_\lambda\|_{L^2(B(x_0,r/2))}^2}{\|\psi_\lambda\|_{L^2(B(x_0,r))}^2}
\ \ge\ \frac{M}{2N_d N}.
\end{equation}

Since $H$ is homogeneous of degree $m$, the rescaling \eqref{eq:rescale} converts \eqref{eq:eigen} into
\begin{equation}\label{eq:scaled-eqn}
Hu = \mu u\quad\text{in }B(0,1),
\quad \mu := r^m\lambda = \frac{\lambda}{|\lambda|}.
\end{equation}
In particular, $|\mu| = 1$.

\textbf{Step 3: find a point where $|u|$ is large.}
By Lemma~\ref{lem:L2-to-Linfty} with $R = 1/2$ and \eqref{eq:u-half}, there exists $y_0\in B(0,1/2)$ such that
\begin{equation}\label{eq:amp-lower}
|u(y_0)|\ \ge\ \frac{\|u\|_{L^2(B(0,1/2))}}{\vol(B(0,1/2))^{1/2}}
\ \ge\ \frac{1}{\vol(B(0,1/2))^{1/2}}\Bigl(\frac{M}{2N_d N}\Bigr)^{1/2}.
\end{equation}

\textbf{Step 4: use a uniform Lipschitz bound to inscribe a nonvanishing ball.}
Since $|\mu| = 1$, Proposition~\ref{prp:uniform-lip} (with $\|u\|_{L^2(B(0,1))} = 1$) yields
$$
|u(x)-u(y)|\le L_{d,H}|x-y|\quad\text{for all }x,y\in B(0,3/4).
$$
Apply Lemma~\ref{lem:lipschitz-ball} to the open set $\Omega' = B(0,3/4)$, the function $G = u|_{\Omega'}$, and the point $x_0 = y_0\in B(0,1/2)\subset\Omega'$.
Since $\dist(y_0,(\Omega')^c) = \frac34-|y_0|\ge \frac14$, Lemma~\ref{lem:lipschitz-ball} gives a ball
$$
B(y_0,\rho_0)\subset\{u\neq 0\},
\quad
\rho_0 := \min\Bigl\{\frac{|u(y_0)|}{2L_{d,H}},\ \frac14\Bigr\}.
$$
Scaling back via $x = x_0+r y$ yields
$$
B(x_0+r y_0,\ r\rho_0)\subset\{\psi_\lambda\neq 0\} = \Sigma_\lambda,
$$
so $\inrad(\Sigma_\lambda)\ge r\rho_0$.

Since $\sqrt{M/N}\le 1$, the elementary inequality $\min\{a t,b\}\ge \min\{a,b\} t$ for $t\in[0,1]$
together with \eqref{eq:amp-lower} yields
$$
\rho_0
\ge
\min\Bigl\{
\frac{1}{2L_{d,H} \vol(B(0,1/2))^{1/2}}\Bigl(\frac{1}{2N_d}\Bigr)^{1/2},
\ \frac14
\Bigr\} \Bigl(\frac{M}{N}\Bigr)^{1/2}.
$$
Therefore
$$
\inrad(\Sigma_\lambda)
\ge
c_{d,H} r \Bigl(\frac{M}{N}\Bigr)^{1/2}
 = 
c_{d,H} r_\lambda 
\frac{\|\psi_\lambda\|_{L^2(\Omega_{-r_\lambda})}}{\|\psi_\lambda\|_{L^2(\Omega)}},
$$
where
$$
c_{d,H} := 
\min\Bigl\{
\frac{1}{2L_{d,H} \vol(B(0,1/2))^{1/2}}\Bigl(\frac{1}{2N_d}\Bigr)^{1/2},
\ \frac14
\Bigr\}>0.
$$
This proves \eqref{eq:main}.
\end{proof}

\end{document}
